% ============================================================
%  Professional Book Template — Revery TeX edition
%
%  A sleek, print-ready book that compiles in the Revery TeX web app
%  (busytex WASM engine) and unchanged on a desktop TeX Live.
%
%  START HERE:  the BOOK METADATA block below is the only part you
%  have to edit to make the book yours.  The demo chapters further
%  down show what the design does — replace them with your own via
%  the \include{chapter/...} lines in MAIN MATTER.
%
%  COMPILE ORDER (a single pass is not enough — the TOC, the
%  bibliography and the index each need their own run):
%      xelatex main  →  bibtex main  →  xelatex main
%                    →  makeindex main  →  xelatex main
%
%  In Revery TeX you do not run these by hand: press Compile and the
%  app reads biblatex's backend and the \makeindex out of the preamble
%  and drives every pass itself.
%
%  THE CONSTRAINTS — why some obvious choices are missing here:
%    • XeLaTeX only        — no LuaLaTeX format ships in the browser engine
%    • No system fonts     — WASM has no host font database.  A font must
%                            come from the TeX distribution itself, so
%                            \setmainfont{Times New Roman} stops the build.
%                            TeX Gyre fonts ship with the engine.
%    • No biber            — biber is a Perl program; only classic BibTeX
%                            (backend = bibtex) exists in the browser engine.
%    • No EPS / PS artwork — that needs Ghostscript; use PNG, JPG or PDF.
%
%  QUICK CUSTOMISATION:
%    • Paper format      → \PaperChoice below   (A4 | B5 | A5 | USTrade | Letter)
%    • Margins           → that size's \GeomOpts in the format dispatch
%    • Text font size    → \RootFontSize below  (10pt | 11pt | 12pt)
%    • Letter spacing    → LetterSpace inside each font block
%    • Line spacing      → uncomment a \setstretch line in TYPOGRAPHY
%    • Accent colour     → \definecolor{AccentColor} in COLOURS
%                          (a black-only interior switch sits right below it)
% ============================================================


% ============================================================
%  BOOK METADATA — edit only these
%  Defined before \documentclass so both the title pages and the PDF
%  document properties (hyperref, further down) can reuse them.
%  Shops, distributors and search engines read that embedded metadata,
%  so it must not be left saying "Book Title".
% ============================================================
\newcommand{\BookTitle}{The Craft of Clear Pages}
\newcommand{\BookSubtitle}{A Practical Introduction to Beautiful Books}
\newcommand{\BookAuthor}{Forename Surname}
\newcommand{\BookYear}{2026}
\newcommand{\BookPublisher}{Revery Press}
\newcommand{\BookISBN}{000-0-000-00000-0}          % leave {} if you have none yet
\newcommand{\BookTags}{book, sample, latex, revery tex}   % tags → PDF metadata

% Optional emblem on the title page.  Put a PNG/JPG/PDF in logo/ and
% point this at it.  A missing file is not an error: the space simply
% stays empty and the page still compiles.
\newcommand{\LogoPath}{logo/logo.png}
% ============================================================


% ---- Root font size (body text) ------------------------------
%  LaTeX's book class supports exactly three sizes:
%      10pt | 11pt | 12pt
\newcommand{\RootFontSize}{11pt}

% ---- Paper format ---------------------------------------------
%  Edit \PaperChoice and nothing else — it drives the class options,
%  the page geometry and the trim size together, so those can never
%  fall out of sync.
%
%    A4       210 × 297 mm  — European standard / review draft  (default)
%    B5       176 × 250 mm  — digest, common trade paperback
%    A5       148 × 210 mm  — compact paperback
%    USTrade  152 × 229 mm  — common US trade paperback
%    Letter   216 × 279 mm  — North-American standard
\newcommand{\PaperChoice}{A4}

% Format dispatch — machinery; no need to edit ------------------
%  Each branch sets two things: the class options and the geometry
%  options (\TrimW/\TrimH are kept for the crop-mark block further
%  down).  To add a new press size, copy one \ifx block and change
%  the \def lines.
%  (Runs before \documentclass, so only TeX primitives are used here.)
%  NOTE: the size names must be defined with \newcommand, exactly
%  like \PaperChoice above — \ifx compares macro meanings, and mixing
%  definition methods makes every comparison silently fail.
\newcommand{\hrvBfive}{B5}
\newcommand{\hrvAfive}{A5}
\newcommand{\hrvUSTrade}{USTrade}
\newcommand{\hrvLetter}{Letter}

% Defaults = A4.
\def\ClassOpts{\RootFontSize, a4paper, twoside, openright}
\def\GeomOpts{a4paper,
  top=25mm, bottom=30mm, inner=25mm, outer=20mm,
  headheight=14pt, headsep=8mm, footskip=12mm}
\def\TrimW{210mm}\def\TrimH{297mm}

\ifx\PaperChoice\hrvBfive
  \def\ClassOpts{\RootFontSize, b5paper, twoside, openright}%
  \def\GeomOpts{b5paper,
    top=22mm, bottom=26mm, inner=22mm, outer=18mm,
    headheight=14pt, headsep=8mm, footskip=12mm}%
  \def\TrimW{176mm}\def\TrimH{250mm}%
\fi
\ifx\PaperChoice\hrvAfive
  \def\ClassOpts{\RootFontSize, a5paper, twoside, openright}%
  \def\GeomOpts{a5paper,
    top=18mm, bottom=22mm, inner=18mm, outer=14mm,
    headheight=13pt, headsep=6mm, footskip=10mm}%
  \def\TrimW{148mm}\def\TrimH{210mm}%
\fi
% US Trade has no named class option — the size comes from geometry.
\ifx\PaperChoice\hrvUSTrade
  \def\ClassOpts{\RootFontSize, twoside, openright}%
  \def\GeomOpts{paperwidth=152mm, paperheight=229mm,
    top=20mm, bottom=24mm, inner=20mm, outer=16mm,
    headheight=14pt, headsep=7mm, footskip=11mm}%
  \def\TrimW{152mm}\def\TrimH{229mm}%
\fi
\ifx\PaperChoice\hrvLetter
  \def\ClassOpts{\RootFontSize, letterpaper, twoside, openright}%
  \def\GeomOpts{letterpaper,
    top=25mm, bottom=30mm, inner=25mm, outer=20mm,
    headheight=14pt, headsep=8mm, footskip=12mm}%
  \def\TrimW{216mm}\def\TrimH{279mm}%
\fi

\expandafter\documentclass\expandafter[\ClassOpts]{book}


% ============================================================
%  FONTS
% ============================================================
\usepackage{fontspec}               % XeLaTeX font selection

% WHY THESE FONTS AND NOT Times New Roman / Arial / Garamond
%  Those faces are proprietary fonts installed by your operating
%  system.  The browser engine has no system font database at all —
%  it can only see fonts that ship with the TeX distribution itself.
%  Naming a system font there does not fall back to something
%  similar; it stops the build.  TeX Gyre is the way out: Pagella is
%  a *metric clone* of Palatino — warm, wide and elegant, a classic
%  book face.  SIL Open Font Licence; ships with TeX Live, so this
%  file also builds unchanged on your desktop.

% --- ACTIVE: TeX Gyre Pagella (Palatino metrics) ---------------
\setmainfont{TeX Gyre Pagella}[     % SIL OFL
  % ---- Text letter spacing ------------------------------------
  %  Units: 1/1000 em.  0 = font default (no change).
  %  Positive → looser, negative → tighter.  Typical range −15…+25.
  LetterSpace = 0,
]
% Headings in a contrasting sans — part of the sleek look.
\setsansfont{TeX Gyre Heros}        % Helvetica metrics
%\setmonofont{TeX Gyre Cursor}      % Courier metrics — uncomment for code

% ── Other fonts that also work in the browser engine ──────────
%  Uncomment ONE block and comment out the ACTIVE block above.

% --- TeX Gyre Termes (Times New Roman metrics; formal, neutral) ---
%\setmainfont{TeX Gyre Termes}
%\setsansfont{TeX Gyre Heros}

% --- TeX Gyre Schola (Century Schoolbook metrics; roomy, warm) ---
%\setmainfont{TeX Gyre Schola}
%\setsansfont{TeX Gyre Heros}

% --- Latin Modern (the classic LaTeX look) ---
%\setmainfont{Latin Modern Roman}
%\setsansfont{Latin Modern Sans}

% ── YOUR OWN FONT FILES — the one way to use anything custom ──
%  Copy the .ttf/.otf beside main.tex and load it *by file*, exactly
%  as below.  Loading by file (Path/Extension) works in the browser;
%  loading by installed family name cannot work there.
%\setmainfont{MyBookFont}[Path=./, Extension=.ttf,
%  BoldFont={MyBookFont-Bold}, ItalicFont={MyBookFont-Italic}]

% ============================================================
%  PAGE GEOMETRY  (print-ready)
%  Margins are chosen by \PaperChoice at the top of this file —
%  there is nothing to uncomment here.  To retune the margins for
%  a size, edit that size's \def\GeomOpts in the format dispatch.
%  Margins follow the rule of thumb: inner > outer, bottom > top.
% ============================================================
\expandafter\usepackage\expandafter[\GeomOpts]{geometry}

%  ── Binding gutter (ask your press) ─────────────────────────
%  Thick books need extra room on the spine side so text does not
%  disappear into the binding.  Rule of thumb: ~5mm under 300 pages,
%  ~8–10mm beyond that.  Adds to the inner margin without moving the
%  trim size, so crop marks stay correct.  Uncomment to enable:
%\geometry{bindingoffset=5mm}

%  ── Trim marks / bleed (ask your press FIRST) ────────────────
%  Default is OFF, and that is what most presses and print-on-demand
%  services want: a plain PDF whose page size IS the trim size.
%  Only switch this on if your press explicitly asked for marks —
%  and note paperwidth/paperheight is the SHEET, not the trim (here:
%  trim + 10mm on every side):
%\usepackage[cam, info=center, center,
%  width=\dimexpr\TrimW+20mm\relax,
%  height=\dimexpr\TrimH+20mm\relax]{crop}
%  For bleed (artwork running off the trimmed edge) ask the press —
%  usually 3mm — and extend images that far past the trim.

% ============================================================
%  LANGUAGE & TYPOGRAPHY
% ============================================================
\usepackage[english]{babel}   % english ships complete; other languages:
%\usepackage{polyglossia}\setmainlanguage{swedish}   % polyglossia works

\usepackage[
  protrusion = true,
  % expansion is not supported by XeTeX — omit entirely
  verbose    = silent,   % suppress microtype slot-number warnings
]{microtype}
\emergencystretch = 1em  % lets TeX relax slightly before an overfull line

% ---- Widow & orphan control -----------------------------------
%  A "widow"  is a paragraph's last  line alone at the TOP of a page;
%  an "orphan" is its first line alone at the BOTTOM.  Both look
%  unprofessional in print.  TeX's maximum penalty tells it to avoid
%  them at virtually any cost.
\widowpenalty = 10000
\clubpenalty  = 10000

%  Books traditionally set \flushbottom (every page exactly the same
%  height); it is the book class default under twoside.  When a page
%  carries mostly unbreakable boxes (tables, listings), stretching
%  becomes impossible and TeX may report an overfull page instead.
%  \raggedbottom lets such pages end short of the margin line:
\raggedbottom

\usepackage{setspace}
% ---- Line spacing ---------------------------------------------
%  \setstretch multiplies the font's natural leading (baseline to
%  baseline).  The default below is the font's own — recommended for
%  a book.  Uncomment ONE value, or set your own number.
%\setstretch{1.0}   % single-spaced — tight; drafts only
%\setstretch{1.15}  % snug — suits technical works
%\setstretch{1.3}   % open — comfortable fiction / large print
%\setstretch{1.5}   % airy — dyslexia-friendly or proofreading copies


% ============================================================
%  COLOURS
% ============================================================
%  Everything coloured in this template routes through these two
%  names: chapter/section headings, caption labels, the title pages
%  and the rules.  The whole interior can therefore be re-skinned
%  here, in one place.
\usepackage[dvipsnames, svgnames, x11names]{xcolor}
% Blue:        1F3A5F
% Green:       2C5F2D
% Brown/Gold:  5E4F2A
% Red:         5E2D30
\definecolor{AccentColor}{HTML}{1F3A5F}   % deep navy — headings & titles
\definecolor{RuleColor}{HTML}{AAAAAA}     % light grey — thin rules

%  ── Black-only interior ─────────────────────────────────────
%  Most text-led books are printed black-only, and that is what
%  print-on-demand quotes assume.  Sent to a black-only press, navy
%  becomes halftone grey, which looks muddy at body size; some
%  presses surcharge or reject colour pages.  Ask your press which
%  interior you are paying for.  For black-only, use these instead:
%\definecolor{AccentColor}{gray}{0}      % pure black headings
%\definecolor{RuleColor}{gray}{0.6}      % light grey rule

% ============================================================
%  CHAPTER & SECTION HEADINGS  (titlesec)
% ============================================================
\usepackage{titlesec}
\newcommand{\headingfont}{\sffamily}   % headings use the sans face

% Chapter style — number above a rule, title below, all sans/navy.
%  \raggedright in both arguments so long titles wrap correctly on
%  narrow pages (A5/USTrade) instead of overflowing the margin.
\titleformat{\chapter}[display]
  {\headingfont\huge\bfseries\color{AccentColor}\raggedright}
  {\chaptertitlename\ \thechapter}
  {8pt}
  {\Huge\raggedright}
  [\vspace{4pt}{\color{RuleColor}\titlerule[0.6pt]}]

% Section style
\titleformat{\section}
  {\headingfont\Large\bfseries\color{AccentColor}}
  {\thesection}{1em}{}

% Subsection style
\titleformat{\subsection}
  {\headingfont\large\bfseries}
  {\thesubsection}{1em}{}

% Slimmer vertical padding around chapter openings than the class default
\titlespacing*{\chapter}{0pt}{-20pt}{28pt}


% ============================================================
%  HEADERS & FOOTERS  (fancyhdr) — correct page numbering
% ============================================================
\usepackage{fancyhdr}
\pagestyle{fancy}
\fancyhf{}
%  Standard book scheme, page numbers on the OUTER edge:
%    verso (even, left-hand)  → page number outer-LEFT, book title inner
%    recto (odd, right-hand)  → chapter title, page number outer-right
\fancyhead[LE]{\small\sffamily\thepage}
\fancyhead[RE]{\small\sffamily\nouppercase{\BookTitle}}
\fancyhead[LO]{\small\sffamily\nouppercase{\leftmark}}
\fancyhead[RO]{\small\sffamily\thepage}
\renewcommand{\headrulewidth}{0.3pt}

% Running head shows just the chapter title, without "Chapter N." —
% the number already appears on the chapter opening page.
\renewcommand{\chaptermark}[1]{\markboth{#1}{}}

% Plain style for chapter-opening pages: bare centred folio, no head.
\fancypagestyle{plain}{%
  \fancyhf{}
  \fancyfoot[C]{\small\sffamily\thepage}
  \renewcommand{\headrulewidth}{0pt}
}


% ============================================================
%  FIGURES, CAPTIONS & PLACEHOLDER ART
% ============================================================
\usepackage{graphicx}
\usepackage{float}
\usepackage[
  labelfont={bf,sf,color=AccentColor},
  textfont=it,
  labelsep=period,
  justification=centering,
  font=small,
]{caption}
\usepackage{subcaption}

% Placeholder boxes when no image file exists yet — drawn with TikZ,
% so the template needs no binary assets to compile.  Replace with
% real art (\includegraphics{graphs/plot1.jpg}) as you go.
\usepackage{tikz}
\newcommand{\placeholder}[3][AccentColor!12]{%
  % Usage: \placeholder[fill colour]{width}{height}
  \begin{tikzpicture}
    \fill[#1] (0,0) rectangle (#2,#3);
    \draw[RuleColor, thick] (0,0) rectangle (#2,#3);
    \node[AccentColor!80] at (#2/2,#3/2) {\small\sffamily\itshape placeholder};
  \end{tikzpicture}%
}


% ============================================================
%  TABLES & CODE
% ============================================================
\usepackage{booktabs}    % professional horizontal rules
\usepackage{tabularx}    % auto-width columns
\usepackage{array}
\usepackage{longtable}   % tables spanning pages
\usepackage{multirow}    % cells spanning multiple rows
\usepackage{wrapfig}     % narrow figures wrapped by running text

% Code listings — styled to match the accent scheme.
\usepackage{listings}
\lstset{
  basicstyle=\ttfamily\small,
  keywordstyle=\color{AccentColor}\bfseries,
  commentstyle=\itshape\color{gray},
  stringstyle=\color{AccentColor},
  numbers=left,
  numberstyle=\tiny\color{RuleColor},
  frame=single,
  rulecolor=\color{RuleColor},
  framesep=6pt,
  aboveskip=12pt,
  belowskip=12pt,
  breaklines=true,
  showstringspaces=false,
}

% ============================================================
%  MISC PACKAGES
% ============================================================
\usepackage{amsmath}     % the modern maths environment set
\usepackage{amssymb}     % blackboard-bold symbols (\mathbb)
\usepackage{amsthm}      % theorem / definition / proof environments

% ---- Theorem-like environments --------------------------------
%  Numbered per chapter ("Theorem 2.1"), sharing one counter so a
%  chapter's results read as one sequence.  Styled to match the
%  book: navy sans header, italic body.
\theoremstyle{definition}
\newtheorem{theorem}{Theorem}[chapter]
\newtheorem{lemma}[theorem]{Lemma}
\newtheorem{definition}[theorem]{Definition}
\newtheorem{example}[theorem]{Example}
\theoremstyle{remark}
\newtheorem*{note}{Note}
\usepackage{enumitem}    % control list spacing
\usepackage{csquotes}    % context-sensitive quotes (required by biblatex)
\usepackage{emptypage}   % keeps inserted blank versos free of heads/numbers

% ---- Index ----------------------------------------------------
%  \makeindex in the preamble is what tells the Revery TeX app to
%  run its makeindex pass.  Use makeidx, NOT imakeidx: imakeidx's
%  automatic mode drives makeindex through shell-escape, which is
%  permanently off in the browser engine.
\usepackage{makeidx}
\makeindex

%  Emptiness test: \IfNonEmpty{<field>}{<code>} typesets <code> only
%  when <field> expands to something non-empty.  This is what lets
%  any metadata field (e.g. \BookISBN) be left empty {} without
%  breaking the compile.
\makeatletter
\newcommand{\IfNonEmpty}[2]{%
  \def\hrv@test{#1}%
  \if\relax\hrv@test\relax
    \expandafter\@gobble
  \else
    \expandafter\@firstofone
  \fi
  {#2}}
\makeatother

% ============================================================
%  BIBLIOGRAPHY  (biblatex on classic BibTeX — never biber)
%  Compile order: xelatex → bibtex → xelatex → xelatex
%
%  WHY bibtex AND NOT biber
%   biber is a Perl program, so no WebAssembly build of it exists.
%   biblatex's older backend, classic BibTeX, *is* compiled into the
%   browser engine (as bibtex8), and biblatex.bst ships with it — so
%   backend = bibtex is the difference between a working bibliography
%   and none at all.  On a desktop TeX Live you may switch to
%   backend = biber for stronger sorting/name handling; delete any
%   stale main.bbl first when changing backends.
%
%  style options:  numeric | authoryear | alphabetic | apa
%  sorting:        nty | nyt | ydnt | none (none = citation order)
% ============================================================
\usepackage[
  backend = bibtex,      % bundled bibtex8; no in-browser engine has biber
  style   = numeric,     % numbered [1] [2] [3] citations
  sorting = none,        % references numbered in order of first citation
]{biblatex}
% Your sources live here (the folder structure this template uses):
\addbibresource{references/reference.bib}

% Hyperlinks & PDF metadata — load LAST.  The properties come from
% the BOOK METADATA block; check them in the finished PDF with pdfinfo.
\usepackage[
  pdfauthor  = {\BookAuthor},
  pdftitle   = {\BookTitle},
  pdfsubject = {\BookSubtitle},
  pdfkeywords= {{\BookTags}},
  pdfcreator = {Revery TeX},
]{hyperref}
\hypersetup{hidelinks}

% ---- cleveref (optional but recommended) ----------------------
%  Turns \cref{label} into "fig. 3", "table 2", "chapter 5" — the
%  word and the number in one command, always matching the target's
%  type.  MUST be loaded AFTER hyperref (it patches hyperref's
%  internals; here it sits last for exactly that reason).
%  cleveref carries its own translations inside cleveref.sty (Swedish
%  included), so it does NOT depend on babel's language files and is
%  safe in the browser engine.  The app's reference picker starts
%  offering \cref automatically once this line is present.
\usepackage[noabbrev]{cleveref}
%\usepackage[noabbrev, swedish]{cleveref}   % Swedish wording instead


% ============================================================
\begin{document}

% ============================================================
%  FRONT MATTER — roman page numbers, recto/verso parity
%  ------------------------------------------------------------
%  Presses expect a fixed hand for each front-matter page:
%    half-title  → recto (right-hand, odd)     title → recto
%    copyright   → verso (back of title leaf)  dedication → recto
%    contents    → recto                       preface → recto
%  \cleardoublepage enforces "start on a recto" by inserting a blank
%  verso when needed; plain \clearpage lands wherever happens to be
%  next.  Use \clearpage ONLY where the next page really is meant to
%  be the immediate verso (title → copyright).  emptypage keeps the
%  inserted blanks clean.
% ============================================================

% ---- Half-title page (recto) ---------------------------------
\frontmatter          % lowercase roman page numbers from here
\thispagestyle{empty}
\vspace*{\fill}
\begin{center}
  {\headingfont\LARGE\bfseries\color{AccentColor}\BookTitle}
\end{center}
\vspace*{\fill}
\cleardoublepage      % blank verso inserted automatically

% ---- Full title page (recto) ---------------------------------
\thispagestyle{empty}
\begin{center}
  \vspace*{12mm}
  % Logo is optional — a missing file simply skips the whole line.
  \IfFileExists{\LogoPath}{\includegraphics[height=25mm]{\LogoPath}\par\vspace{14mm}}{}%
  {\headingfont\Huge\bfseries\color{AccentColor}\BookTitle}\par
  \vspace{6mm}
  {\color{RuleColor}\rule{0.5\textwidth}{0.6pt}}\par
  \vspace{7mm}
  {\Large\itshape\BookSubtitle}\par
  \vfill
  {\large\BookAuthor}\par
  \vspace{10mm}
  {\headingfont\small\color{AccentColor}\BookPublisher}\par
  \vspace{2mm}
  {\small\BookYear}\par
  \vspace*{12mm}
\end{center}
\clearpage            % copyright is the immediate verso of the title leaf

% ---- Copyright page (verso) ----------------------------------
\thispagestyle{empty}
\vspace*{\fill}
\begin{small}
  \noindent
  Copyright \textcopyright{} \BookYear\ \BookAuthor\par
  \vspace{6pt}
  All rights reserved. No part of this publication may be reproduced
  or transmitted in any form without written permission of the
  author.\par
  \vspace{6pt}
  \IfNonEmpty{\BookISBN}{ISBN \BookISBN\par\vspace{6pt}}%
  Typeset with \LaTeX{} using the Revery TeX book template.\par
\end{small}
\vspace*{\fill}
\cleardoublepage

% ---- Dedication (recto, optional) -----------------------------
\thispagestyle{empty}
\vspace*{5cm}
\begin{center}
  \itshape To everyone who reads to the last page.
\end{center}
\cleardoublepage

% ---- Table of contents (recto) --------------------------------
\tableofcontents
\cleardoublepage


% ---- Preface ---------------------------------------------------
%  RULE: every \chapter* needs its own TOC line and running marks.
%  \chapter* does NOT set the running marks nor a contents entry —
%  without the two lines below, the running head would keep saying
%  "Contents" and the preface would be missing from the TOC.  Repeat
%  this pattern for Foreword, Acknowledgements, About the Author…
\chapter*{Preface}
\addcontentsline{toc}{chapter}{Preface}
\markboth{Preface}{}

This book was typeset with \LaTeX{} using the Revery TeX book
template. Replace this sample text with your own preface — or delete
the whole section if your book does not need one.\index{preface}


% ============================================================
%  MAIN MATTER — arabic page numbers restart at 1 here
% ============================================================
\mainmatter

% ---- Demo chapters — replace with your own --------------------
% ============================================================
%  chapter/intro.tex — Getting started
%
%  This chapter doubles as a mini-manual.  Write your real first
%  chapter in this file, or add new files in chapter/ and pull them
%  in from main.tex with \include{chapter/yourfile}.
% ============================================================
\chapter{Getting Started}\label{ch:intro}

Welcome to your book. Everything you see on these pages is sample
content that exists to show off the design and to stress-test the
compiler. Replace it freely.

\section{Writing text}
Ordinary paragraphs need no markup at all: just write, and leave one
blank line between paragraphs. \emph{Italics}, \textbf{bold} and
\textsc{small caps} are inline commands; footnotes go at the bottom of
the page.\footnote{Like this one. Footnotes are numbered per chapter.}

Quotation marks are best typed with the \texttt{\textbackslash enquote} command from
\emph{csquotes}, which adapts to the document language:
\enquote{like this}, rather than by typing quote characters by hand.

\subsection{Lists}
Three kinds come built in, tuned by the \emph{enumitem} package:
\begin{itemize}[itemsep=2pt]
  \item an itemised list,
  \item like this one;
  \item nesting works to any sane depth.
\end{itemize}
\begin{enumerate}[itemsep=2pt]
  \item an enumerated list,
  \item numbered automatically;
  \item cross-referenced anywhere as (\ref{enum:start}).
        \label{enum:start}
\end{enumerate}

\section{Cross-references}
Every chapter, section, figure, table and equation can carry a
\texttt{\textbackslash label}; refer to it with \texttt{\textbackslash ref} or
\texttt{\textbackslash eqref}. This chapter is
Chapter~\ref{ch:intro}, and the maths chapter is
Chapter~\ref{ch:maths}. Compile twice so the numbers settle.

The template also loads \emph{cleveref}, which adds the type word
for you: \cref{enum:start} needs no typed ``item'', and
\cref{ch:maths} comes out as a full ``chapter'' reference. Use
\texttt{\textbackslash cref} mid-sentence and \texttt{\textbackslash Cref} to open one.

\section{Citations and the index}
Cite anything in \texttt{references/reference.bib} by its key: the
classic advice is Knuth~\cite{knuth1984}, and modern practice is well
summarised elsewhere~\cite{lamport1994, mittal2005}. The full list is
printed automatically in the back matter, numbered in order of first
citation.

For the index, wrap any word in \texttt{\textbackslash index}: this book talks about
\index{typesetting}typesetting, \index{page breaks}page breaks and
the \index{index}index itself. Sub-entries use an exclamation mark:
\index{index!sub-entries}like so.

\section{A little longer prose}\label{sec:prose}
Good book typography is invisible: the reader should feel the page,
not notice it.\footnote{The classic reference here is
\textcite{bringhurst2004}, whose rules on measure, leading and rag
still hold on screen as on paper.} Three settings do most of that
work in this template --- measure (the line length set by the paper
format), leading (the line spacing in TYPOGRAPHY) and the colour of
the text block (how evenly the words grey the page).

The measure deserves respect. A line of running text wants between
sixty and seventy-five characters; shorter tires the eye with turns,
longer loses the return journey to the left margin. The margins in
each \texttt{\textbackslash GeomOpts} block are tuned with
that in mind, which is why they differ per paper size.

\section{Description lists}
A third list kind, for term--definition pairs:

\begin{description}[itemsep=2pt]
  \item[Measure] the width of the text block, usually given in ems or
    characters per line.
  \item[Leading] the baseline-to-baseline distance; the template's
    default is the font's own.
  \item[Folio] the printed page number, set on the outer edge here.
\end{description}

\noindent
The sources behind these terms, and everything else in this chapter,
are collected in the bibliography: see \parencite{knuth1984,
lamport1994} for the typesetting machinery and \textcite{texgyre}
for the fonts themselves.\index{leading}\index{measure|see{line length}}

% ============================================================
%  chapter/chapter1.tex — Mathematics stress test
% ============================================================
\chapter{Mathematics}\label{ch:maths}

Inline maths such as $e^{i\pi} + 1 = 0$ sits in the running text;
display maths gets its own line and number:

\begin{equation}\label{eq:gauss}
  \int_{-\infty}^{\infty} e^{-x^{2}} \,\mathrm{d}x = \sqrt{\pi}.
\end{equation}

Equation~\eqref{eq:gauss} is referenced with \texttt{\textbackslash eqref}. Multi-line
derivations use \texttt{align}:

\begin{align}
  (a+b)^2 &= a^2 + 2ab + b^2,\\
  (a-b)^2 &= a^2 - 2ab + b^2,\\
  (a-b)(a+b) &= a^2 - b^2.
\end{align}

Matrices and cases work through \texttt{amsmath}:

\begin{equation}
  A = \begin{pmatrix}
        a_{11} & a_{12} & \cdots & a_{1n}\\
        a_{21} & a_{22} & \cdots & a_{2n}\\
        \vdots & \vdots & \ddots & \vdots\\
        a_{m1} & a_{m2} & \cdots & a_{mn}
      \end{pmatrix},
  \qquad
  f(x) = \begin{cases}
           x^{2}, & x \geq 0,\\
           -x,    & x < 0.
         \end{cases}
\end{equation}

Theorem-like statements can be built with the classic
\texttt{amsthm} machinery if your book needs it; for most trade
books the styled paragraph below is enough.

\section{Symbols in the index}
Technical terms should reach the index as you introduce them: the
Euler identity\index{Euler identity}, Gaussian
integrals\index{Gaussian integral} and
matrices\index{matrix|see{linear algebra}} all appear above.

\section{Statements and proofs}
Formal books state, then prove. The environments come from
\emph{amsthm}, styled in this template's navy-and-sans scheme:

\begin{definition}\label{def:prime}
A natural number $p > 1$ is \emph{prime} if its only positive
divisors are $1$ and $p$ itself.
\end{definition}

\begin{theorem}[Euclid]\label{thm:primes}
There are infinitely many primes.
\end{theorem}

\begin{proof}
Suppose $p_1, p_2, \dots, p_n$ were all the primes and consider
$N = p_1 p_2 \cdots p_n + 1$. No $p_i$ divides $N$, since each
leaves remainder~$1$; yet $N > 1$ has a prime factor by
Definition~\ref{def:prime}. Contradiction.
\end{proof}

\begin{example}
The first few primes are $2, 3, 5, 7, 11, 13$;
by Theorem~\ref{thm:primes} the list never ends.\index{prime number}
\end{example}

\begin{note}
Unnumbered remarks use the starred \texttt{note} environment ---
useful for asides that would clutter the numbering.
\end{note}

\section{More display mathematics}
Fractions, binomials and named operators behave as usual:
\begin{equation}\label{eq:binom}
  \binom{n}{k} = \frac{n!}{k!\,(n-k)!},
  \qquad \sum_{k=0}^{n} \binom{n}{k} = 2^{n},
  \qquad \lim_{x \to 0} \frac{\sin x}{x} = 1.
\end{equation}
Text belongs inside maths through \texttt{\textbackslash text},
which keeps the surrounding font:
\begin{equation*}
  f(x) = x^{2} \quad \text{for all } x \in \mathbb{R}.
\end{equation*}
Both numbered forms cross-reference: combine
\eqref{eq:gauss} with \eqref{eq:binom} and you have most of a
first calculus course.\index{binomial coefficient}

% ============================================================
%  chapter/chapter2.tex — Figures & tables stress test
%
%  Real artwork goes into graphs/ and illustrations/.  Until then
%  the \placeholder boxes show where images will sit.
% ============================================================
\chapter{Figures and Tables}\label{ch:figures}

\section{Figures}
Drop your plots into \texttt{graphs/} and load them with
\texttt{\textbackslash includegraphics}; the placeholder below shows the frame the
finished figure will occupy.

\begin{figure}[ht]
  \centering
  % Real plot: \includegraphics[width=0.8\linewidth]{graphs/plot1.jpg}
  \placeholder[AccentColor!12]{0.8\linewidth}{55mm}
  \caption{A placeholder for a data plot. Replace the
    \texttt{\textbackslash placeholder} line with an
    \texttt{\textbackslash includegraphics} pointing at
    \texttt{graphs/plot1.jpg}.}
  \label{fig:plot1}
\end{figure}

Figure~\ref{fig:plot1} floats to the top of a page when that helps
the layout; \texttt{[ht]} lets LaTeX choose between ``here'' and
``top''. Side-by-side sub-figures work too:

\begin{figure}[ht]
  \centering
  \begin{subfigure}{0.48\linewidth}
    \centering
    \placeholder{0.9\linewidth}{35mm}
    \caption{First panel.}
  \end{subfigure}\hfill
  \begin{subfigure}{0.48\linewidth}
    \centering
    \placeholder{0.9\linewidth}{35mm}
    \caption{Second panel.}
  \end{subfigure}
  \caption{Two sub-figures sharing one number.}
  \label{fig:panels}
\end{figure}

\section{Tables}
Booktabs rules (no vertical lines) are the professional standard:

\begin{table}[ht]
  \centering
  \caption{A small comparison table.}
  \label{tab:compare}
  \begin{tabular}{lcc}
    \toprule
    Format & Trim size & Typical use\\
    \midrule
    A4       & 210 × 297 mm & review drafts\\
    B5       & 176 × 250 mm & digest editions\\
    US Trade & 152 × 229 mm & trade paperback\\
    \bottomrule
  \end{tabular}
\end{table}

Table~\ref{tab:compare} uses fixed columns; for text-heavy tables
that must fit the measure exactly, use \texttt{tabularx}:

\begin{table}[ht]
  \centering
  \caption{A tabularx table — the X column absorbs the spare width.}
  \label{tab:x}
  \begin{tabularx}{\linewidth}{lX}
    \toprule
    Term & Meaning\\
    \midrule
    Verso & The left-hand, even-numbered page of a spread.\\
    Recto & The right-hand, odd-numbered page of a spread.\\
    Folio & The page number itself, printed on the outer edge here.\\
    \bottomrule
  \end{tabularx}
\end{table}

\section{A wrapped side figure}
Narrow figures can sit beside the running text instead of floating
away from it --- useful for portraits and small diagrams:

\begin{wrapfigure}{r}{0.42\linewidth}
  \centering
  % Real portrait: \includegraphics{illustrations/portrait.jpg}
  \placeholder{0.9\linewidth}{34mm}
  \caption{A figure wrapped by the text.}
  \label{fig:wrap}
\end{wrapfigure}

Figure~\ref{fig:wrap} occupies the right of the measure while these
paragraphs flow around it. Keep wrapped figures narrow --- about a
third to two-fifths of the text width --- so the remaining lines stay
long enough to read comfortably. Two or three paragraphs is the most
text that should share a column with a wrapfigure; beyond that the
page starts to look pinched. Note how the wrapping has now finished:
a section heading should never have to share its lines with a
floating column, so leave enough prose after the figure to close the
wrap before the next heading arrives.

\section{Spanning tables and grouped rows}
Two table problems come up constantly: cells that span several rows,
and tables longer than a page. The first is \texttt{multirow}'s job,
shown in Table~\ref{tab:grouped} --- which, unlike the tables so
far, is \emph{not} a float: it is set exactly here with
\texttt{\textbackslash captionof}, so it can never drift away from
the sentence that introduces it. Use this for small tables whose
position matters; let everything else float.

% A non-floating table: \captionof (from the caption package) keeps
% the numbering and the List of Tables entry without any floating.
\begin{center}
  \captionof{table}{Rows grouped under spanning cells.}
  \label{tab:grouped}
  \begin{tabular}{llr}
    \toprule
    Family & Face & Weights\\
    \midrule
    \multirow{3}{*}{TeX Gyre} & Pagella & regular, bold, italic\\
                              & Heros   & regular, bold\\
                              & Termes  & regular, bold, italic\\
    \addlinespace
    \multirow{2}{*}{Latin Modern} & Roman & regular, bold\\
                                  & Sans  & regular, bold\\
    \bottomrule
  \end{tabular}
\end{center}

The second problem --- a table that outgrows the page --- belongs to
\texttt{longtable}, which breaks across pages and repeats its header
on every continuation:

\begin{longtable}{@{}p{28mm}p{\dimexpr\linewidth-28mm-12pt\relax}@{}}
  \caption{A long table that spans pages.}\label{tab:long}\\
  \toprule
  Term & Definition\\
  \midrule
  \endfirsthead
  \multicolumn{2}{@{}l}{\small\itshape Table~\ref{tab:long}, continued}\\
  \toprule
  Term & Definition\\
  \midrule
  \endhead
  \bottomrule
  \endfoot
  Ascender   & The part of a letter that rises above the x-height, as in b, d, k.\\
  Baseline   & The invisible line on which letters appear to sit.\\
  Descender  & The part that drops below the baseline, as in p, q, y.\\
  Kerning    & Selective adjustment of space between specific letter pairs.\\
  Ligature   & Two or more letters cast as one glyph, such as fi or fl.\\
  Serif      & The small finishing stroke at the ends of a letter's main strokes.\\
  Terminal   & The often-rounded end of a stroke, as in the tail of an a.\\
  X-height   & The height of the lowercase letters, excluding ascenders.\\
\end{longtable}

% ============================================================
%  chapter/code.tex — Source-code listings stress test
%
%  The listing style is configured in main.tex (LISTINGS section).
% ============================================================
\chapter{Code Listings}\label{ch:code}

Technical books often carry source code. The \emph{listings} package
typesets it with line numbers, a frame and keyword colouring that
match the book's accent scheme:

\begin{lstlisting}[language=Python, caption={A tiny Python function.}, label={lst:py}]
def fibonacci(n):
    """Return the n-th Fibonacci number."""
    a, b = 0, 1
    for _ in range(n):
        a, b = b, a + b
    return a

print([fibonacci(i) for i in range(10)])
\end{lstlisting}

Listing~\ref{lst:py} floats like any figure when given the float
option; inline snippets need no environment at all:
\lstinline|printf("%d\n", x);|.

\begin{lstlisting}[language=C, caption={C, for good measure.}, label={lst:c}]
#include <stdio.h>

int main(void) {
    /* Hello from the listings package */
    for (int i = 0; i < 3; i++)
        printf("line %d\n", i);
    return 0;
}
\end{lstlisting}

Long listings can break across pages (\texttt{breaklines=true} is
already set), and captions place them in the optional List of
Listings if you ever add \texttt{\textbackslash lstlistoflistings} to the front matter.
Without a language, the block typesets verbatim --- right for
config files and shell transcripts:

\begin{lstlisting}[caption={A language-free (verbatim) listing.}, label={lst:sh}]
# build the book from a clean folder
rm -f main.aux main.toc main.bbl main.idx
xelatex main && bibtex main && xelatex main
makeindex main && xelatex main
\end{lstlisting}

Structured data reads well too:

\begin{lstlisting}[caption={A JSON fragment.}, label={lst:json}]
{
  "title": "The Craft of Clear Pages",
  "year": 2026,
  "keywords": ["book", "sample", "latex"]
}
\end{lstlisting}

Listing~\ref{lst:sh} drives the full compile; listing~\ref{lst:json}
shows that braces survive intact inside any listing.

% ============================================================
%  chapter/examples.tex — Text patterns & citation showcase
%
%  A grab-bag of the devices books actually use: quotations, verse,
%  epigraphs, footnote-dense prose, and every common way to cite.
% ============================================================
\chapter{Text Patterns}\label{ch:patterns}

\section{Quotations}
Short quotations stay inline \enquote{like this one}; anything worth
its own lines uses the quotation environment:

\begin{quotation}
\noindent\itshape
Typography exists to honour content.\index{typography!purpose}
\end{quotation}
\begin{flushright}
--- \citeauthor{bringhurst2004}, \citeyear{bringhurst2004}
\end{flushright}

A displayed extract from another work keeps its original lineation
in the \texttt{verse} environment:

\begin{verse}
To see a World in a Grain of Sand\\
And a Heaven in a Wild Flower\\
\hspace*{2em}Hold Infinity in the palm of your hand\\
\hspace*{2em}And Eternity in an hour
\end{verse}

\section{Epigraphs}
Chapter openings often carry a short attributed line above the
first heading. There is no package for that here --- just a styled
block before the first section:

\begin{quote}
\small\itshape\color{AccentColor}
The form of a book must be answerable to the energy of its contents.
\end{quote}

\section{Citations, every common flavour}\label{sec:cites}
Biblatex offers one command per rhetorical situation, all reading
the same \texttt{reference.bib}. In running prose, name the source:
\textcite{knuth1984} designed TeX;
\textcite{lamport1994} put LaTeX on top of it. Parenthetical
asides take \parencite{mittal2005}; multiple sources collapse into
one bracket \parencite{knuth1984,lamport1994,mittal2005}.

When the year matters, write \textcite{texgyre}
(\citeyear{texgyre}); when the title itself is quoted,
\emph{\citetitle{knuth1984}} does the job. For a note at the bottom
of the page, use a footnote citation\footcite{bringhurst2004} ---
the full entry prints below the line without interrupting the
paragraph, a favourite in humanities and history
titles.\index{citations}

Every cited work lands automatically in the back-matter
bibliography, numbered in order of first appearance because the
template sets \texttt{sorting=none}.\index{bibliography}

\section{Footnotes in numbers}
Dense scholarly prose leans on notes.\footnote{Like this. Notes are
per-chapter here, resetting with each \texttt{\textbackslash chapter}.} They
should carry evidence,\footnote{Or the source, as footnotes do.}
qualification,\footnote{Or a polite disagreement.} or a joke too
good to cut\footnote{This one, arguably.} --- but never the argument
itself. If a reader must open your note to follow the sentence, the
sentence is in the wrong place.

\section{An index-rich close}
Everything a reader might hunt for belongs in the index, nested
where it helps:\index{book design}\index{book design!margins}%
\index{book design!grids}\index{page layout|seealso{book design}}
this closing section demonstrates sub-entries under a shared head,
and a \emph{see also} redirect. Compile, check
\texttt{main.ind}, and trust nothing you have not seen on a page.

% ---- Appendix -------------------------------------------------
\appendix            % letters the chapter numbers from here on
% ============================================================
%  chapter/appendix.tex — Appendix
%
%  Loaded after \appendix in main.tex, so this chapter is numbered
%  "Appendix A" automatically.
% ============================================================
\chapter{Notation}\label{app:notation}

Appendices hold reference material: notation tables, worked proofs,
glossaries, long source listings. A compact notation table makes a
good first appendix:

\begin{center}
  \begin{tabular}{ll}
    \toprule
    Symbol & Meaning\\
    \midrule
    $\mathbb{N}$   & natural numbers\\
    $\mathbb{Z}$   & integers\\
    $\mathbb{R}$   & real numbers\\
    $\in$          & element of\\
    $\subseteq$    & subset of (possibly equal)\\
    $\sum_{k=1}^{n} a_k$ & the sum $a_1 + a_2 + \cdots + a_n$\\
    \bottomrule
  \end{tabular}
\end{center}

\noindent
Terms that readers may hunt for later belong in the index as well as
in the text:\index{notation}\index{appendix}.



% ============================================================
%  BACK MATTER
% ============================================================
\backmatter

% ---- Bibliography ----------------------------------------------
%  The sources live in references/reference.bib.  biblatex adds the
%  heading and the TOC entry itself; heading=bibintoc puts
%  "Bibliography" in the contents.
\printbibliography[heading=bibintoc, title={Bibliography}]

% ---- Index -----------------------------------------------------
%  Add \index{term} commands anywhere in the text.  \printindex
%  prints the heading but no TOC entry, so add one by hand.
\cleardoublepage
\addcontentsline{toc}{chapter}{Index}
\printindex


% ---- Signature padding (ask your press) ------------------------
%  Books are printed on folded sheets, so many presses require the
%  final page count to be a multiple of 4 (sometimes 8 or 16).  If
%  needed, add blank leaves with lines like these until it fits;
%  emptypage keeps them free of headers and numbers:
%\clearpage\thispagestyle{empty}\null\clearpage   % repeat as needed


% ============================================================
%  PRE-FLIGHT CHECKLIST — run through before sending to press
% ============================================================
%  1. DEMO CONTENT OUT.  Remove or replace the sample chapters above,
%     the placeholder figures and the sample citations once your own
%     text is in.
%  2. FONTS EMBEDDED.      pdffonts main.pdf   → every row "emb yes"
%  3. NO OVERFULL LINES.   grep -c "Overfull" main.log  → expect 0
%  4. PAGE COUNT multiple of 4/8/16 — see signature padding above.
%  5. IMAGES at 300 dpi at their final printed size; PNG/JPG/PDF only.
%  6. TRIM MARKS / BLEED off unless the press asked for them.
%  7. INTERIOR INK: black-only vs colour — set COLOURS accordingly.
%  8. METADATA correct:    pdfinfo main.pdf
%  9. FULL CLEAN BUILD: delete main.aux/.toc/.bbl/.idx and
%     chapter/*.aux first, then xelatex → bibtex → xelatex →
%     makeindex → xelatex.

\end{document}




